\documentstyle[%


righttag%


,11pt%


,amssymb%


,russian%               remove in Eng.


,izvru%                 Set izven% in Eng.


]{amsart}





% special attn line 550


\newcommand\supp{\operatorname{supp}}


\newcommand\ind{\operatorname{ind}}


\newcommand\esssup{\operatornamewithlimits{ess\,sup}}


\newcommand\mes{\operatorname{mes}}


\newcommand{\dst}{\displaystyle}


\newcommand{\const}{\operatorname{const}}


\newcommand{\RE}{\operatorname{Re}}


\newcommand{\Div}{\operatorname{Div}}


\newcommand{\Hom}{\operatorname{Hom}}


\newcommand{\trivert}{|@!@!|@!@!|} % adjust spacing, p.9


\newcommand{\apA}{\overset{\;\approx}{A}{}} % adjust


\newcommand{\apsigma}{\overset{\approx}{\sigma}{}}


\newcommand{\xO}{\overset{x_0}{\sim}{}}


\newcommand{\tts}[1]{}





%\hoffset=-15mm%                  For Eng.


\setcounter{page}{3}


\god{2001}


\nomer{1~(464)} %                        Remove in Eng.


%\nomer{Vol.\,45, No.\,1}%               Set in Eng.


%\str{\pageref{begin}--\pageref{end}}%  Set in Eng.


\udk{517.9}





\author{\it О.Г.\,Авсянкин, Н.К.\,Карапетянц}


% NB:   ^^ remove in Eng.


\title{Об алгебре многомерных интегральных операторов


с однородными ядрами с переменными коэффициентами}


\thanks{Работа поддержана Российским фондом


фундаментальных исследований, проект \No\,98-01-00261а.}





\date{}


%\pagestyle{myheadings}%         Set in Eng.


\pagestyle{plain} %              Remove in Eng.


\begin{document}


%\thispagestyle{plain}%          Set in Eng.


\maketitle


\label{begin}





Хорошо известна роль, которую играют интегральные операторы с


однородными ядрами в различных разделах математики и в приложениях.


Изучение таких операторов в многомерной ситуации было начато в


работах Л.Г.\,Михайлова и его учеников (см. \cite{b1}, а также


библиографию в \cite{b1}) и


продолжено в \cite{b2}, \cite{b3}. В настоящее время для


интегральных операторов с ядрами, однородными степени $(-n)$ и


инвариантными относительно всех вращений в $\Bbb R^n\times\Bbb R^n$,


построена полная н\"етеровская теория, теория проекционных методов, а


также исследованы банаховы алгебры, порожденные такими операторами.





В данной работе рассматривается вопрос о н\"етеровости операторов из


банаховой алгебры $\frak B$, порожденной многомерными интегральными


операторами с однородными ядрами и операторами умножения на функции


из класса $\frak A$, действующими в $L_p(\Bbb R^n)$, $1<p<\infty$


(точная постановка задачи и определение класса $\frak A$ будут даны


ниже). На алгебре $\frak B$ строится символическое исчисление, в


терминах которого описываются необходимые и достаточные условия


н\"етеровости операторов из этой алгебры и указывается формула для


вычисления индекса.


Отметим, что основным средством решения указанной задачи является


локальный метод И.Б.\,Симоненко \cite{b4}.





В статье используются следующие обозначения:


$\Bbb R^n$ --- $n$-мерное евклидово пространство; $x=(x_1,\dotsc,


x_n)\in\Bbb R^n$; $|x|=\sqrt{x_1^2+\dotsb +x_n^2}$; $x'=x/|x|$;


$x\cdot y=x_1y_1+\dotsb+x_ny_n$; $dx=dx_1\dotsc dx_n$; $e_1=(1,0,\dotsc,


0)$; $\Dot{\Bbb R}^n$ --- компактификация $\Bbb R^n$ одной бесконечно


удаленной точкой; $y=\omega_{x'}(t)$ --- вращение в $\Bbb R^n$,


переводящее $t\in\Bbb R^n$ в $y\in\Bbb R^n$ так, что


$\omega_{x'}(e_1)=x'$; $\Bbb Z_+$ --- множество всех целых


неотрицательных чисел; $\Bbb Z_+\dot\times\Bbb R$ ---


компактификация множества $\Bbb Z_+\times\Bbb R$ одной бесконечно


удаленной точкой; $P_m(t)$ --- многочлены Лежандра, определяемые


следующим образом:


$$


P_m(t)=


\begin{cases}


\cos(m\arccos t), &n=2;\\


(C_{m+n-3}^m)^{-1}C_m^{(n-2)/2}(t), &n\geqslant 3,


\end{cases}


$$


где $C_m^{(n-2)/2}(t)$ --- многочлены Гегенбауэра.





\section{О компактности интегральных операторов с однородными ядрами с


переменными коэффициентами}





В пространстве $L_p(\Bbb R^n)$, $1<p<\infty$,


рассмотрим интегральный оператор


\begin{equation}


(K\varphi)(x)=\int_{\Bbb R^n} k(x,y)\varphi(y)\,dy,\quad


x\in\Bbb R^n,


\label{1}


\end{equation}


предполагая, что ядро $k(x,y)$ удовлетворяет условиям


\begin{itemize}


\item[$1^\circ$.] однородности степени $(-n)$, т.\,е.


$k(\lambda x,\lambda y)=\lambda^{-n}k(x,y)$


$\forall\, \lambda>0$;


\item[$2^\circ$.] инвариантности относительно группы $SO(n)$ вращений


пространства $\Bbb R^n$, т.\,е.


$k(\omega(x),\linebreak \omega(y))=k(x,y)$


$\forall\, \omega\in SO(n)$;


\item[$3^\circ$.] суммируемости, т.\,е.


$k=\dst\int_{\Bbb R^n}|k(e_1,y)||y|^{-n/p}dy<\infty$.


\end{itemize}


Функцию $a(x)$ будем естественно отождествлять с оператором


умножения на эту функцию.





Нас будут интересовать условия компактности оператора $a(x)K$.


Известно \cite{b2}, что если $a(x)\in L_\infty^{0,


0}(\Bbb R^n)$, то оператор $a(x)K$ компактен в $L_p(\Bbb R^n)$.


Укажем более слабые условия на коэффициент $a(x)$,


обеспечивающие компактность оператора $a(x)K$.





Обозначим через $\frak A_0$ множество всех функций $a(x)\in


L_\infty(\Bbb R^n)$ таких, что для любого компактного


множества $M$ выполняются условия


\begin{align}


\lim_{\alpha\to 0}\int_M |a(\alpha x)|dx&=0,\label{2.2}\\


\lim_{\alpha\to \infty}\int_M |a(\alpha x)|dx&=0.


\label{2.3}


\end{align}





\begin{lem}\label{l2.1}


Пусть $a(x)\in\frak A_0$. Тогда для


любой функции $f(x)\in L_1(\Bbb R^n)$ выполняются условия


\begin{align}


\lim_{\alpha\to 0}\int_{\Bbb R^n} |a(\alpha x)f(x)|dx&=0,


\label{2.4}\\


\lim_{\alpha\to \infty}\int_{\Bbb R^n} |a(\alpha x)f(x)|


dx&=0.\label{2.5}


\end{align}


\end{lem}





\begin{pf}


Докажем (\ref{2.4}). 


Возьмем любое $\varepsilon>0$. Подберем функцию $g(x)\in


C_0^\infty(\Bbb R^n)$ такую, что


$$


\int_{\Bbb R^n} |f(x)-g(x)|dx<


\frac{\varepsilon}{2\|a\|_\infty},


$$


и зафиксируем функцию $g(x)$. Имеем


\begin{multline*}


\int_{\Bbb R^n} |a(\alpha x)f(x)|dx\leqslant


\int_{\Bbb R^n} |a(\alpha x)|\,|f(x)-g(x)|dx+


\int_{\Bbb R^n} |a(\alpha x)|\,|g(x)|dx\leqslant\\


%


\leqslant \|a\|_\infty \int_{\Bbb R^n} |f(x)-g(x)|dx+


\int_{\supp g}|a(\alpha x)|\,|g(x)|dx<


\frac{\varepsilon}{2}+\max_{x\in\Bbb R^n}|g(x)|


\int_{\supp g}|a(\alpha x)|dx.


\end{multline*}


В силу (\ref{2.2}) найдется такое $\alpha_0>0$, что при любом


$\alpha\in(0,\alpha_0)$ выполняется условие


$$


\int_{\supp g}|a(\alpha x)|dx<


\frac{\varepsilon}{2\max\limits_{x\in\Bbb R^n}|g(x)|}\,.


$$


Таким образом, неравенство


$\dst\int_{\Bbb R^n}|a(\alpha x)f(x)|\,dx<\varepsilon$


справедливо для всех $\alpha\in(0,\alpha_0)$, что и доказывает


(\ref{2.4}). Условие (\ref{2.5}) доказывается аналогично.


\end{pf}





Обозначим через $K_a$ оператор $Ka(x)I$.





\begin{lem}\label{l2.2}


Для оператора $K_a$


\begin{equation}


\|K_a\|\leqslant k^{1/p}\|a\|_\infty^{1/p}


\esssup_{x\in\Bbb R^n}


\bigg(\int_{\Bbb R^n}


|k(e_1,t)|\,|t|^{-n/p}


\big|a(|x|\omega_{x'}(t))\big|dt\bigg)^{1/p'}.


\label{2.6}


\end{equation}


\end{lem}





{\bf Доказательство.}


Воспользовавшись неравенством Г\"ельдера и затем выполняя вращение


$y=\omega_{x'}(t)$, получаем


\begin{multline*}


|(K_a\varphi)(x)|\leqslant


\int_{\Bbb R^n}|k(x,y)|\,|a(y)|\,|\varphi(y)|dy \leqslant\\


%


\leqslant\bigg(\int_{\Bbb R^n}|k(x,y)|\,|y|^{-n/p}


|a(y)|^{p'}dy\bigg)^{1/p'}\bigg(\int_{\Bbb R^n}|k(x,


y)|\,|y|^{n/p'}|\varphi(y)|^pdy\bigg)^{1/p}=\\


%


= |x|^{-\frac{n}{pp'}}\bigg(\int_{\Bbb R^n}|k(e_1,


t)|\,|t|^{-n/p}|a(|x|\omega_{x'}(t))|^{p'}dt\bigg)^{1/p'}


\bigg(\int_{\Bbb R^n}|k(x,


y)|\,|y|^{n/p'}|\varphi(y)|^p dy\bigg)^{1/p}.


\end{multline*}


Далее,


\begin{multline*}


\|K_a\varphi\|_p^p\leqslant


\esssup_{x\in\Bbb R^n}


\bigg(\int_{\Bbb R^n}|k(e_1,t)|\,|t|^{-n/p}


\big|a(|x|\omega_{x'}(t))\big|^{p'}dt\bigg)^{p/p'}\times\\


%


\times\int_{\Bbb R^n}|x|^{-n/p'}dx


\int_{\Bbb R^n}|k(x,y)|\,|y|^{n/p'}|\varphi(y)|^p dy\leqslant\\


%


\leqslant\|a\|_\infty^{(p'-1)p/p'}


\esssup_{x\in\Bbb R^n}


\bigg(\int_{\Bbb R^n}|k(e_1,t)|\,|t|^{-n/p}


|a(|x|\omega_{x'}(t))|dt\bigg)^{p/p'}\times\\


%


\times\int_{\Bbb R^n}|\varphi(y)|^p|y|^{n/p'}dy


\int_{\Bbb R^n}|k(x,y)|\,|x|^{-n/p'}dx=\\


%


=k\|a\|_\infty^{(p'-1)p/p'}


\|\varphi\|_p^p\esssup_{x\in\Bbb R^n}


\bigg(\int_{\Bbb R^n}|k(e_1,t)|\,|t|^{-n/p}


|a(|x|\omega_{x'}(t))|dt\bigg)^{p/p'}.


\end{multline*}


Отсюда следует неравенство


$$


\|K_a\varphi\|_p\leqslant k^{1/p}\|a\|_\infty^{1/p}


\esssup_{x\in\Bbb R^n}


\bigg(\int_{\Bbb R^n}|k(e_1,t)|\,|t|^{-n/p}


|a(|x|\omega_{x'}(t))|dt\bigg)^{1/p'}\|\varphi\|_p.


\quad\square


$$





\begin{thm}\label{t2.1}


Пусть $a(x)\in\frak A_0$. Тогда


оператор $K_a$ компактен в $L_p(\Bbb R^n)$, $1<p<\infty$.


\end{thm}





\begin{pf}


Так как


$|k(e_1,t)|\,|t|^{-n/p}\in L_1(\Bbb R^n)$, то в силу леммы~1


$$


\int_{\Bbb R^n}|k(e_1,t)|\,|t|^{-n/p}


|a(|x|\omega_{x'}(t))|dt\to 0\quad


\text{при}\quad|x|\to0


$$


и


$$


\int_{\Bbb R^n}|k(e_1,t)||t|^{-n/p}


|a(|x|\omega_{x'}(t))|dt


\to 0\quad


\text{при}\quad|x|\to\infty.


$$


Следовательно, для любого $\varepsilon >0$ найдется такое число


$N\in \Bbb N$, что


\begin{equation}


\esssup_{|x|<\frac1N\cup |x|>N}


\int_{\Bbb R^n}|k(e_1,t)|\,|t|^{-n/p}


|a(|x|\omega_{x'}(t))|dt<\varepsilon.


\label{2.7}


\end{equation}


Обозначим $D_N=\{x\in\Bbb R^n: 1/N<|x|<N\}$ и пусть $P_{D_N}$


--- оператор умножения на характеристическую функцию множества


$D_N$. Известно (см. \cite{b2}), что оператор $P_{D_N}K_a$


компактен. Далее, с помощью леммы~2, имеем


\begin{multline*}


\|K_a-P_{D_N}K_a\| =


\|P_{\Bbb R^n\setminus D_N}K_a\|\leqslant\\


%


\leqslant k^{1/p}\|a\|_\infty^{1/p}


\esssup_{x\in \Bbb R^n\setminus D_N}


\bigg(\int_{\Bbb R^n}|k(e_1,t)|\,|t|^{-n/p}


|a(|x|\omega_{x'}(t))|dt\bigg)^{1/p'}<


k^{1/p} \|a\|_\infty^{1/p} \varepsilon^{1/p'}


\end{multline*}


в силу (\ref{2.7}). Итак, оператор $K_a$ можно приблизить


компактными операторами $P_{D_N}K_a$ с любой степенью точности.


Значит, оператор $K_a$ компактен.


\end{pf}





Из теоремы~1 посредством перехода к сопряженному оператору


получается





\begin{thm}


Пусть $a(x)\in\frak A_0$.


Тогда оператор $a(x)K$ компактен в $L_p(\Bbb R^n)$,


$1<p<\infty$.


\end{thm}





В заключение отметим, что для операторов свертки аналоги условий


(\ref{2.2}) и (\ref{2.3}) были сформулированы в \cite{b5}, \cite{b6}.





\section{Определение класса $\frak A$}





Обозначим через $\frak A$ множество всех функций $a(x)\in


L_\infty(\Bbb R^n)$, для каждой из которых существуют числа $a_0$ и


$a_\infty$ такие, что для любого компактного множества


$M$ выполняются условия


\begin{align*}


\lim_{\alpha\to 0}\int_M |a(\alpha x)-a_0|dx&=0,\\


\lim_{\alpha\to \infty}\int_M |a(\alpha x)-a_\infty|dx&=0.


\end{align*}


Очевидно, $\frak A_0=\{a(x)\in\frak A: a_0=a_\infty=0\}$.





Далее, обозначим через $\chi_+(x)$ и $\chi_-(x)$


характеристические функции единичного шара в $\Bbb


R^n$ и его внешности соответственно. Каждой функции


$a(x)\in\frak A$ сопоставим кусочно-постоянную функцию


$\wt a(x)=a_0\chi_+(x)+a_\infty\chi_-(x)$.





\begin{lem}


Пусть $a(x)\in\frak A$. Тогда функция


$b(x)=a(x)-\wt a(x)$ принадлежит классу $\frak A_0$.


\end{lem}





\begin{pf}


Для любого компакта $M$ имеем


$$


\lim_{\alpha\to0} \int_M |b(\alpha x)|dx=


\lim_{\alpha\to0} \int_M


|a(\alpha x)-a_0\chi_+(\alpha x)-a_\infty\chi_-(\alpha x)|dx=


\lim_{\alpha\to0} \int_M |a(\alpha x)-a_0|dx=0.


$$


Таким образом, условие (\ref{2.2}) выполнено.





Проверим условие (\ref{2.3}).


Вначале возьмем любой компакт $M_0$, не содержащий нуля.


Тогда получим


$$


\lim_{\alpha\to\infty} \int_{M_0} |b(\alpha x)|dx=


\lim_{\alpha\to\infty} \int_{M_0}


|a(\alpha x)-a_0\chi_+(\alpha x)-a_\infty\chi_-(\alpha x)|dx=


\lim_{\alpha\to\infty} \int_{M_0}


|a(\alpha x)-a_\infty|dx=0.


$$


Пусть теперь $M$ --- произвольный компакт в $\Bbb R^n$. Возьмем


произвольно $\varepsilon>0$ и подберем компакт $M_0\subset M$ такой,


что $0\notin M_0$ и $\mes(M\setminus


M_0)<\varepsilon/2\|b\|_\infty$. Тогда


$$


\int_M |b(\alpha x)|dx= \int_{M_0}


|b(\alpha x)|dx+\int_{M\setminus M_0}|b(\alpha x)|dx


< \int_{M_0} |b(\alpha x)|dx+\frac{\varepsilon}{2}.


$$


В силу сказанного выше первое слагаемое стремится к нулю при


$\alpha\to\infty$. Следовательно, найдется такое число $\alpha_0$,


что для всякого $\alpha>\alpha_0$ справедливо неравенство


$$


\int_M |b(\alpha x)|dx<\varepsilon,


$$


т.\,е. выполняется (\ref{2.3}).


\end{pf}





Из этой леммы и теоремы 2 сразу следует равенство


$$


I-a(x)K=I-\wt a(x)K+T,


$$


где $T=-b(x)K$ --- компактный в $L_p(\Bbb R^n)$, $1<p<\infty$,


оператор. Соответственно


\begin{equation}


\sum_i\prod_j(I-a_{ij}(x)K_{ij})=


\sum_i\prod_j(I-\wt a_{ij}(x)K_{ij})+T_1,


\label{9}


\end{equation}


где $T_1$ --- компактный оператор, а сумма и произведение


предполагаются конечными.





\section{Алгебра $\frak B$}





Обозначим через $\frak B$ наименьшую замкнутую подалгебру банаховой


алгебры $\cal L(L_p(\dot{\Bbb R}^n))$, $1<p<\infty$, содержащую все


операторы вида $I-a(x)K$, где $K$ --- оператор вида (\ref{1}), а


$a(x)\in\frak A$. Она представляет собой замыкание по операторной


норме множества


$$


\frak B_0=\bigg\{\sum_i\prod_j(I-a_{ij}(x)K_{ij})\bigg\},


$$


где сумма и произведение конечны.





Наша цель состоит в построении на алгебре $\frak B$ символического


исчисления, в терминах которого описываются необходимые и


достаточные условия н\"етеровости операторов из этой алгебры.





Напомним (см. \cite{b4}), что оператор $A:L_p(\Dot{\Bbb R}^n)\to


L_p(\Dot{\Bbb R}^n)$ называется оператором локального типа, если для


любой пары замкнутых непересекающихся множеств $F_1$ и $F_2$


оператор $P_{F_1}AP_{F_2}$ компактен.





\begin{lem}


Пусть $A\in\frak B$. Тогда оператор $A$ является оператором


локального типа.


\end{lem}





\begin{pf}


Вначале докажем, что оператор $K$


является оператором локального типа.


Пусть $F_1$ и $F_2$ --- замкнутые непересекающиеся множества в


$\Dot{\Bbb R}^n$, $\chi_1(x)$ и $\chi_2(x)$ --- характеристические


функции этих множеств. Рассмотрим оператор


$$


(P_{F_1}KP_{F_2}\varphi)(x)=


\int_{\Bbb R^n}\chi_1(x)\chi_2(y)k(x,y)\varphi(y)dy,\quad


x\in\Dot{\Bbb R}^n.


$$


На основании результатов \cite{b2} заключаем, что оператор


$P_{F_1}KP_{F_2}$ компактен в $L_p(\Dot{\Bbb R}^n)$, $1<p<\infty$.


Следовательно, $K$ --- оператор локального типа.





Известно \cite{b4}, что оператор умножения на существенно


ограниченную функцию является оператором локального типа. Поэтому


оператор $I-a(x)K$, где $a(x)\in\frak A$, есть оператор локального


типа. Но тогда и любой оператор $A\in\frak B_0$ является оператором


локального типа (\cite{b4}, с.\,569).





Пусть теперь $A$ --- произвольный оператор из алгебры $\frak B$.


Тогда найдется такая последовательность $\{A_k\}\subset\frak B_0$,


что $\|A_k-A\|\to 0$ при $k\to\infty$. Тогда тем более


$\trivert A_k-A\trivert\to0$ при $k\to\infty$. (Здесь и далее


$\trivert C\trivert=\inf\limits_T\|C+T\|$, где $T$ пробегает множество всех


компактных операторов.) Учитывая свойство операторов локального типа


\cite{b4}, заключаем, что $A$ является оператором


локального типа.


\end{pf}





Используем два понятия, введенные в


\cite{b4}. Пусть $A$ и $B$ --- операторы локального типа. Операторы


$A$ и $B$ называются эквивалентными в точке $x_0$, если для любого


$\varepsilon>0$ найдется такая окрестность $u$ точки $x_0$, что


$\trivert(A-B)P_u\trivert<\varepsilon$ (обозначение $A\xO B$).


Локальным представителем оператора $A$ в точке $x_0$ будем называть


любой оператор, эквивалентный оператору $A$ в этой точке.





\begin{lem}


Для того чтобы оператор


\begin{equation}


A=\sum_i\prod_j (I-a_{ij}(x)K_{ij})\in\frak B_0


\label{10}


\end{equation}


являлся н\"етеровым, необходимо и достаточно, чтобы н\"етеровыми были


операторы


\begin{equation}


\wt A=\sum_i\prod_j (I-(a_0)_{ij}K_{ij})\quad


\text{и}\quad


\apA=\sum_i\prod_j (I-(a_\infty)_{ij}K_{ij}).


\label{11}


\end{equation}


\end{lem}





{\bf Доказательство} разобьем на два этапа.





1. Рассмотрим оператор $I-a(x)K$, где $a(x)\in\frak A$, а $K$ ---


оператор вида (\ref{1}). Построим локальные представители


оператора


$I-a(x)K$ в каждой точке $x_0\in\Dot{\Bbb R}^n$.





1) Пусть $x_0\ne 0$ и $x_0\ne\infty$. Тогда, очевидно, что


$I-a(x)K \xO I$.





2) Пусть $x_0=0$. Покажем, что


$I-a(x)K \xO I-a_0K$.





Действительно, если $u$ --- некоторая окрестность точки $x_0=0$, то


оператор $[a(x)-a_0]KP_u$ является компактным в $L_p(\dot{\Bbb


R}^n)$ оператором. Следовательно,


$\trivert[(I-a(x)K)-(I-a_0K)]P_u


\trivert=\trivert[a(x)-a_0]KP_u\trivert=0$,


что и требовалось.





3) Пусть $x_0=\infty$. Рассуждая аналогично предыдущему, получаем


$I-a(x)K \xO I-a_\infty K$.





2. Из вышесказанного с учетом свойств операторов локального


типа (см. \cite{b4}) следует





а) $A\xO I$, если $x_0\ne 0$ и $x_0\ne \infty$,





б) $A\xO \wt A$, если $x_0=0$,





в) $A\xO \apA$, если $x_0=\infty$.





\noindent{}Тогда в силу теоремы 1.6 из \cite{b4}


оператор $A$ н\"етеров тогда и


только тогда, когда н\"етеровы операторы $\wt A$ и $\apA$.$\quad\square$





Назовем символом оператора $I-a(x)K$ пару функций


$\big(\wt\sigma(m,\xi),\apsigma(m,\xi)\big)$, определенных на компакте


$\Bbb Z_+\Dot\times\Bbb R$ равенствами


\begin{align*}


\wt\sigma(m,\xi)&=1-a_0\int_{\Bbb R^n}


k(e_1,y)P_m(e_1\cdot y')|y|^{-n/p+i\xi}dy,\\


%


\apsigma(m,\xi)&=1-a_\infty\int_{\Bbb R^n}


k(e_1,y)P_m(e_1\cdot y')|y|^{-n/p+i\xi}dy.


\end{align*}


Далее, если $\big(\wt\sigma_{ij}(m,\xi),\apsigma_{ij}(m,\xi)\big)$


--- символ оператора $I-a_{ij}(x)K_{ij}$, то символом оператора $A$,


заданного равенством (\ref{10}), назовем пару функций


$\big(\wt\sigma(m,\xi),\apsigma(m,\xi)\big)$, где


$\wt\sigma(m,\xi)=\sum\limits_i\prod\limits_j \wt\sigma_{ij}(m,\xi)$,


$\apsigma(m,\xi)=\sum\limits_i\prod\limits_j \apsigma_{ij}(m,\xi)$.





\begin{lem}


Для того чтобы оператор $A$, заданный равенством $(\ref{10})$, являлся


н\"етеровым, необходимо и


достаточно, чтобы его символ был невырожденным, т.\,е.


\begin{equation}


\wt\sigma(m,\xi)\ne 0,\quad \apsigma(m,\xi)\ne 0\quad


\forall\, (m,\xi)\in \Bbb Z_+\Dot\times\Bbb R.


\label{12}


\end{equation}


\end{lem}





\begin{pf}


Согласно результатам, полученным в


\cite{b3}, условия (\ref{12}) являются


необходимыми и достаточными условиями н\"етеровости операторов


$\wt A$ и $\apA$ соответственно. И остается лишь применить


лемму~5.


\end{pf}





Установим далее условия н\"етеровости произвольного оператора из


алгебры $\frak B$.





\begin{lem}


Пусть $A$ --- оператор вида $(\ref{10})$. Тогда справедливы неравенства


\begin{equation}


\|\wt A\|\leqslant \trivert A \trivert, \quad


\|\apA\|\leqslant\trivert A\trivert.


\label{13}


\end{equation}


\end{lem}





\begin{pf}


Докажем первое из неравенств (\ref{13}).


Обозначим через $P_+$ оператор умножения на функцию $\chi_+(x)$ и


положим $\cal K_+=P_+KP_+$. На основании результатов \cite{b3}


заключаем, что


$\sum\limits_i\prod\limits_j (I-(a_0)_{ij}K_{ij})=I-aK$,


где $a\in\Bbb C$ и $K$ --- оператор вида (\ref{1}). Тогда


$\sum\limits_i\prod\limits_j (P_+-(a_0)_{ij}\cal K_{ij})=P_+-a\cal K+T$,


где $\cal K=P_+KP_+$, $\cal K_{ij}=P_+K_{ij}P_+$, а $T$ --- компактный


оператор (см. \cite{b3}).





С учетом (\ref{9}) имеем


\begin{multline*}


\trivert A\trivert\ =


\Big|@!@!\Big|@!@!\Big|\sum_i\prod_j (I-a_{ij}(x)K_{ij})


\Big|@!@!\Big|@!@!\Big|=


\Big|@!@!\Big|@!@!\Big| \sum_i\prod_j(I-\wt a_{ij}(x)K_{ij})


\Big|@!@!\Big|@!@!\Big| \geqslant \\


%


\geqslant\Big|@!@!\Big|@!@!\Big| P_+\Big(\sum_i\prod_j(I-\wt


a_{ij}(x)K_{ij})\Big)P_+\Big|@!@!\Big|@!@!\Big| =


\Big|@!@!\Big|@!@!\Big| \sum_i\prod_j P_+(I-\wt a_{ij}(x)K_{ij})P_+


\Big|@!@!\Big|@!@!\Big| =\\


%


= \Big|@!@!\Big|@!@!\Big|\sum_i\prod_j(P_+-(a_0)_{ij}\cal K_{ij})


\Big|@!@!\Big|@!@!\Big|=


\trivert P_+-a\cal K+T\trivert\ =\trivert P_+-a\cal K\trivert =


\trivert P_+(I-aK)P_+\trivert.


\end{multline*}


Учитывая, что


$\trivert P_+(I-aK)P_+\trivert=\|I-aK \|$ (см. \cite{b3}), получаем


$\trivert A\trivert\geqslant \|I-aK\|=


\Big\|\sum\limits_i\prod\limits_j(I-(a_0)_{ij}K_{ij})\Big\|=\|\wt A \|$.


\end{pf}





Теперь определим понятие символа для произвольного оператора из


алгебры $\frak B$. Если $A\in\frak B_0$, то это сделано выше. Пусть


$A\in\frak B\setminus\frak B_0$. Тогда найдется такая


последовательность $\{A_k\}\subset\frak B_0$, что $\|A_k-A\|\to 0$


при $k\to\infty$. Тогда тем более $\trivert A_k-A\trivert\to 0$ при


$k\to\infty$. Каждому оператору $A_k\in\frak B_0$ вида (\ref{10})


соответствует пара операторов $\wt A_k$ и $\apA_k$ вида (\ref{11}).


Рассмотрим последовательности операторов $\{\wt A_k\}$ и


$\{\apA_k\}$. Так как последовательность $\{A_k\}$ фундаментальна по


существенной норме, то в силу (\ref{13}) последовательности


$\{\wt A_k\}$ и $\{\apA_k\}$ фундаментальны по операторной норме.


Положим $\dst\wt A=u-\lim_{k\to\infty}\wt A_k$ и


$\dst \apA=u-\lim_{k\to\infty} \apA_k$. Обозначим через


$\wt\sigma(m,\xi)$ и $\apsigma(m,\xi)$ символы операторов $\wt A$


и $\apA$ соответственно. Пару функций


($\wt\sigma(m,\xi)$,\ $\apsigma(m,\xi)$) будем называть символом


оператора $A$.





\begin{thm}


Для того чтобы оператор $A$ являлся н\"етеровым, необходимо и


достаточно, чтобы его символ был невырожденным, т.\,е.


\begin{equation}


\wt\sigma(m,\xi)\ne 0,\quad


\apsigma(m,\xi)\ne 0\quad \forall\,(m,\xi)\in \Bbb Z_+\Dot\times\Bbb R.


\label{14}


\end{equation}


Если условия $(\ref{14})$ выполнены, то индекс оператора $A$ вычисляется


по формуле


\begin{equation}


\ind A=\sum_{m=0}^\infty d_n(m)\ind_\xi


(\apsigma(m,\xi)/\wt\sigma(m,\xi)), \label{15}


\end{equation}


где $\ind_\xi f(m,\xi)$ --- индекс функции $f(m,\xi)$ при


фиксированном значении $m$.


\end{thm}





\begin{pf}


Если $A\in \frak B_0$, то теорема уже


доказана (см. лемму~6). Пусть


$A\in\frak B\setminus\frak B_0$. Рассмотрим последовательность


$\{A_k\}\subset\frak B_0$ такую, что $\|A_k-A\|\to 0$ при


$k\to\infty$. Положим $\dst \wt A=u-\lim_{k\to\infty}\wt A_k$ и


$\dst \apA=u-\lim_{k\to\infty}\apA_k$.


Построим локальные представители оператора $A$ в каждой точке


$x_0\in\Bbb R^n$.





1) Пусть $x_0\ne 0$ и $x_0\ne\infty$. Так как


$\trivert A_k-A\trivert\to 0$ при $k\to\infty$ и $A_k\xO I$, то $A\xO I$.





2) Пусть $x_0=0$. Так как $\trivert A_k-A\trivert\to 0$ при


$k\to\infty$, $\trivert \wt A_k-\wt A\trivert\to 0$ при


$k\to\infty$ и $A_k\xO\wt A_k$, то $A\xO\wt A$.





3) Пусть ${x_0=\infty}$. Так как ${\trivert A_k-A\trivert\to 0}$ при


${k\to\infty}$,


$\trivert \apA_k-\apA\trivert\to 0$ при $k\to\infty$ и


$A_k\xO\apA_k$, то $A\xO\apA$.





Условия (\ref{14}) являются необходимыми и достаточными


условиями н\"етеровости операторов $\wt A$ и $\apA$. Тогда в силу


теоремы 1.6 из \cite{b4} условия (\ref{14}) являются необходимыми и


достаточными для н\"етеровости оператора $A$.





Установим теперь формулу (\ref{15}). Если оператор $A$ н\"етеров, то


существует такое $\varepsilon>0$, что для всех операторов $B$


таких, что $\|B\|<\varepsilon$, оператор $A+B$ н\"етеров и


$\ind(A+B)=\ind A$.





Так как $A\overset{0}{\sim}{}\wt A$ и


$A\overset{\infty}{\sim}{}\apA$, то по выбранному $\varepsilon$


найдутся окрестности $u_1$ точки $x=0$ и $u_2$ точки $x=\infty$


такие, что


$$


\trivert (A-\wt A)P_{u_1}\trivert <\frac\varepsilon 2\quad


\text{и}\quad


\trivert (A-\apA)P_{u_2}\trivert <\frac\varepsilon 2.


$$


Тогда тем более


$$


\trivert P_{u_1}(A-\wt A)P_{u_1}\trivert <\frac\varepsilon 2\quad


\text{и}\quad


\trivert P_{u_2}(A-\apA)P_{u_2}\trivert <\frac\varepsilon 2.


$$


В силу определения существенной нормы найдутся такие компактные


операторы $T_1$ и $T_2$, что будут выполнены неравенства


\begin{equation}


\|P_{u_1}(A-\wt A)P_{u_1}+T_1\| <\frac\varepsilon 2\quad


\text{и}\quad


\| P_{u_2}(A-\apA)P_{u_2}+T_2\| <\frac\varepsilon 2.


\label{16}


\end{equation}


Полагая $P=I-P_{u_1}-P_{u_2}$, получим


$$


A=(P_{u_1}+P+P_{u_2})A(P_{u_1}+P+P_{u_2})=


P_{u_1}AP_{u_1}+P_{u_2}AP_{u_2}+T,


$$


где $T$ --- некоторый компактный оператор. Представим оператор $A$ в


виде


$$


A=P_{u_1}\wt AP_{u_1} +P_{u_2}\apA P_{u_2} +(T-T_1-T_2)+B,


$$


где


$$


B=P_{u_1}(A-\wt A)P_{u_1} +T_1+P_{u_2}(A-\apA)P_{u_2} +T_2.


$$


Из неравенств (\ref{16}) следует, что $\|B\|<\varepsilon$. Тогда


$\ind(A-B)=\ind A$.


Поскольку $A-B=P_{u_1}\wt AP_{u_1} +P_{u_2}\apA P_{u_2}


+T-T_1-T_2$, то


$$


\ind(A-B)=\ind(P_{u_1}\wt AP_{u_1} +P_{u_2}\apA P_{u_2})=


\ind P_{u_1}\wt AP_{u_1} +\ind P_{u_2}\apA P_{u_2}.


$$


Таким образом,


$\ind A=\ind P_{u_1}\wt AP_{u_1} +\ind P_{u_2}\apA P_{u_2}$.


Из результатов, полученных в \cite{b3}, следует, что


\begin{align*}


\ind P_{u_1}\wt AP_{u_1}&=


-\sum_{m=0}^\infty d_n(m)\ind_\xi \wt\sigma(m,\xi),\\


\ind P_{u_2}\apA P_{u_2}&=


\sum_{m=0}^\infty d_n(m)\ind_\xi \apsigma(m,\xi).


\end{align*}


Отсюда следует формула (\ref{15}).


\end{pf}





{\bf Замечание.} Следует отметить, что поскольку, начиная с


некоторого $m_0$, все индексы\linebreak $\ind_\xi (\apsigma(m,


\xi)/\wt\sigma(m,\xi))$ равны нулю, то правая часть в формуле


(\ref{15}) фактически представляет собой конечную сумму.





\begin{thebibliography}{9}





\bibitem{b1}


Михайлов Л.Г. {\em Новый класс особых интегральных уравнений} // Math.


Nachr. -- 1977. -- Т.\,76. -- С.\,91--107.





\bibitem{b2}


Карапетянц Н.К. {\em Интегральные операторы свертки и с однородными


ядрами с переменными коэффициентами}: Дисс.~$\dotsc$ докт. физ.-матем.


наук. -- Тбилиси, 1989. -- 296~с.





\bibitem{b3}


Авсянкин О.Г. {\em Многомерные интегральные операторы с однородными


ядрами}: Дисс.~$\dotsc$ канд. физ.-матем. наук. --


Ростов-на-Дону, 1997. -- 145~с.





\bibitem{b4}


Симоненко И.Б. {\em Новый общий метод исследования линейных операторных


уравнений типа сингулярных интегральных уравнений}. I // Изв. АН СССР. Сер.


матем. -- 1965. -- Т.\,29. -- \No\,3. -- С.\,567--586.





\bibitem{b5}


Speck~F.O. {\em Eine Erweiterung des Satzes von Rakov\v c{}ik und ihre


Anwendung in der Simonenko--Theorie} // Math. Ann. --


1977. -- V.\,228. -- \No\,2. -- P.\,93--100.





\bibitem{b6}


Штейнберг Б.Я. {\em Об операторах типа свертки на локально компактных


группах} // Функц. анализ и его прилож. -- 1981. --


Т.\,15. -- Вып.\,3. -- С.\,95--96.








\end{thebibliography}





\vskip 0.5cm





\post{\it Ростовский государственный университет}{\it Поступила}


\post{}{05.02.1999}





\label{end}


\end{document}


